% FAC-EN-U0061
% Source: fac_body.tex lines 2270--2339.
% Authority: PDF physical pp. 61--62; printed pp. 256--257
% (Chapter III, no. 63 complete).
% Printed-source correction C0037 restores the superscript 0 on L_n^0.

\phantomsection\label{fac:III:3:no:63}
\textbf{63. General properties of $H^q(M)$.}

\phantomsection\label{fac:III:3:no:63:prop:3}
\textsc{Proposition 3.} \textit{Let $M$ be a graded $S$-module
\hyperref[fac:III:2:no:56:def:TF]{satisfying condition}}
\hyperref[fac:III:2:no:56:def:TF]{(TF)}. \textit{Then:}

\phantomsection\label{fac:III:3:no:63:prop:3:item:a}
\textit{(a) There exists an integer $k(M)$ such that
$\rho_k^h:H_k^q(M)\to H_h^q(M)$ is bijective for
$h\geq k\geq k(M)$ and every $q$.}

\phantomsection\label{fac:III:3:no:63:prop:3:item:b}
\textit{(b) $H^q(M)$ is a finite-dimensional vector space over $K$ for
every $q\geq0$.}

\phantomsection\label{fac:III:3:no:63:prop:3:item:c}
\textit{(c) There exists an integer $n(M)$ such that, for $n\geq n(M)$,
$\alpha:M_n\to H^0(M(n))$ is bijective and $H^q(M(n))=0$ for every
$q>0$.}

We reduce at once to the case in which $M$ is of finite type. We shall
then say that $M$ has \emph{dimension} $\leq s$ (where $s$ is an integer
$\geq0$) if there is an exact sequence
\phantomsection\label{fac:III:3:no:63:def:dimension}
\[
0\longrightarrow L^s\longrightarrow L^{s-1}\longrightarrow\cdots
\longrightarrow L^0\longrightarrow M\longrightarrow0,
\]
in which the $L^i$ are graded free $S$-modules of finite type. By
Hilbert's syzygy theorem (see [6], Chapter VIII, Theorem 6.5), this
dimension is always $\leq r+1$.

We prove \hyperref[fac:III:3:no:63:prop:3]{the proposition} by induction
on the dimension of $M$. If that
dimension is zero, $M$ is free of finite type, that is, a direct sum of
modules $S(n_i)$, and the proposition follows from
\hyperref[fac:III:3:no:62:cor:2]{Corollaries~2} and
\hyperref[fac:III:3:no:62:cor:3]{3} to
\hyperref[fac:III:3:no:62:prop:2]{Proposition~2}. Suppose that $M$
has dimension $\leq s$, and let $N$ be the kernel of $L^0\to M$. The
graded $S$-module $N$ has dimension $\leq s-1$, and there is an exact
sequence
\phantomsection\label{fac:III:3:no:63:formula:short-exact}
\[
0\longrightarrow N\longrightarrow L^0\longrightarrow M
\longrightarrow0.
\]
By the induction hypothesis,
\hyperref[fac:III:3:no:63:prop:3]{the proposition} holds for $N$ and $L^0$.
Applying the five lemma ([7], Chapter I, Lemma 4.3) to the commutative
diagram
\phantomsection\label{fac:III:3:no:63:diagram:five-lemma}
\[
\begin{tikzcd}
H_k^q(N) \arrow[r] \arrow[d] & H_k^q(L^0) \arrow[r] \arrow[d]
  & H_k^q(M) \arrow[r] \arrow[d] & H_k^{q+1}(N) \arrow[r] \arrow[d]
  & H_k^{q+1}(L^0) \arrow[d] \\
H_h^q(N) \arrow[r] & H_h^q(L^0) \arrow[r] & H_h^q(M) \arrow[r]
  & H_h^{q+1}(N) \arrow[r] & H_h^{q+1}(L^0),
\end{tikzcd}
\]
for $h\geq k\geq\operatorname{Sup}(k(N),k(L^0))$, proves
\hyperref[fac:III:3:no:63:prop:3:item:a]{(a)}, and hence also
\hyperref[fac:III:3:no:63:prop:3:item:b]{(b)}, since the $H_k^q(M)$ are
finite-dimensional over $K$.
On the other hand, the exact sequence
\phantomsection\label{fac:III:3:no:63:formula:cohomology-segment}
\[
H^q(L^0(n))\longrightarrow H^q(M(n))\longrightarrow H^{q+1}(N(n))
\]
shows that $H^q(M(n))=0$ for
$n\geq\operatorname{Sup}(n(L^0),n(N))$. Finally, consider the
commutative diagram
\phantomsection\label{fac:III:3:no:63:diagram:alpha}
\[
\begin{tikzcd}
0 \arrow[r] \arrow[d] & N_n \arrow[r] \arrow[d,"\alpha"]
  & L_n^0 \arrow[r] \arrow[d,"\alpha"]
  & M_n \arrow[r] \arrow[d,"\alpha"] & 0 \arrow[d] \\
0 \arrow[r] & H^0(N(n)) \arrow[r] & H^0(L^0(n)) \arrow[r]
  & H^0(M(n)) \arrow[r] & H^1(N(n));
\end{tikzcd}
\]
The printed upper middle term is $L_n$, with its superscript $0$
omitted. For $n\geq n(N)$ we have $H^1(N(n))=0$; it follows that
$\alpha:M_n\to H^0(M(n))$ is bijective for
$n\geq\operatorname{Sup}(n(L^0),n(N))$, which completes the proof of
\hyperref[fac:III:3:no:63:prop:3]{the proposition}.
